\documentclass[12pt,a4paper]{article}
\usepackage[UTF8]{ctex}
\usepackage{amsmath,amssymb,amsthm}
\usepackage{geometry}
\usepackage{tikz}
\usepackage{pgfplots}
\usetikzlibrary{shapes,arrows,positioning,calc,decorations.markings}
\geometry{left=2.5cm,right=2.5cm,top=2.5cm,bottom=2.5cm}

\title{平面力闭合：四个Wrench和两个摩擦接触详解}
\author{}
\date{}

\begin{document}

\maketitle

\section{引言}

\subsection{核心陈述}

\textbf{原文}：
\begin{quote}
对于平面问题，四个接触wrench足以正张成三维wrench空间，这意味着至少两个摩擦接触（每个有两条摩擦锥边）足以实现力闭合。使用力矩标记，我们看到力闭合等价于没有一致的力矩标记。例如，如果两个接触可以通过位于两个摩擦锥内的视线"看到"彼此，则我们有力闭合（图12.22(b)）。
\end{quote}

\subsection{需要理解的关键点}

\begin{enumerate}
\item 为什么四个wrench足以正张成三维wrench空间？
\item 为什么两个摩擦接触（每个有两条边）等于四个wrench？
\item 为什么力闭合等价于没有一致的力矩标记？
\item 什么是"视线"（line of sight）？为什么它能判断力闭合？
\end{enumerate}

\section{第一部分：四个Wrench足以正张成三维空间}

\subsection{平面Wrench空间的维度}

\textbf{平面问题}：
\begin{itemize}
\item 物体在2D平面中运动
\item Wrench空间是\textbf{3维}的：$(m_z, f_x, f_y)$
   \begin{itemize}
   \item $m_z$：绕$z$轴的力矩（1维）
   \item $f_x, f_y$：平面内的力（2维）
   \end{itemize}
\end{itemize}

\subsection{正张成需要多少个向量？}

\textbf{线性代数事实}：
\begin{itemize}
\item 要\textbf{线性张成}（linear span）$n$维空间，需要至少$n$个向量
\item 要\textbf{正张成}（positive span）$n$维空间，需要至少$n+1$个向量
\end{itemize}

\textbf{原因}：
\begin{itemize}
\item 线性张成允许负系数：$\sum k_i v_i$，$k_i \in \mathbb{R}$
\item 正张成只允许非负系数：$\sum k_i v_i$，$k_i \geq 0$
\item 正张成更受限制，所以需要更多向量
\end{itemize}

\textbf{对于平面wrench空间}：
\begin{itemize}
\item 维度：$n = 3$
\item 正张成需要：$n + 1 = 4$个向量
\item 因此：\textbf{四个wrench足以正张成三维wrench空间}
\end{itemize}

\subsection{为什么不能更少？}

\textbf{三个wrench的情况}：

假设三个wrench $F_1, F_2, F_3$线性无关（可以线性张成整个空间）。

\textbf{问题}：它们能正张成整个空间吗？

\textbf{答案}：\textbf{不能}！

\textbf{原因}：
\begin{itemize}
\item 对于任何三个向量，存在一个方向$c$，使得对所有$i$有$F_i^T c \leq 0$
\item 这意味着：三个向量的正线性组合不能创建方向$c$的向量
\item 因此不能正张成整个空间
\end{itemize}

\textbf{例子}：
\begin{itemize}
\item $F_1 = (1, 0, 0)$，$F_2 = (0, 1, 0)$，$F_3 = (0, 0, 1)$
\item 正张成：第一象限（只能表示$m_z \geq 0, f_x \geq 0, f_y \geq 0$的点）
\item 不能表示负方向的wrench
\end{itemize}

\textbf{需要第四个wrench}：
\begin{itemize}
\item 例如：$F_4 = (-1, -1, -1)$
\item 现在四个wrench可以正张成整个空间
\end{itemize}

\section{第二部分：两个摩擦接触 = 四个Wrench}

\subsection{摩擦锥的表示}

\textbf{平面摩擦锥}：
\begin{itemize}
\item 对于平面问题，摩擦锥是\textbf{2维}的（在接触点的切平面内）
\item 摩擦锥由\textbf{两条边}定义（见图12.18(b)）
\item 每条边对应一个\textbf{方向}的摩擦力
\end{itemize}

\textbf{摩擦锥的两条边}：
\begin{itemize}
\item 边1：摩擦力方向与接触法线成角度$+\alpha$（$\alpha = \tan^{-1}\mu$）
\item 边2：摩擦力方向与接触法线成角度$-\alpha$
\item 两条边张成整个摩擦锥
\end{itemize}

\subsection{从摩擦锥边到Wrench}

\textbf{每个摩擦锥边对应一个wrench}：
\begin{itemize}
\item 接触位置：$p_i$
\item 摩擦锥边方向：$\hat{d}_i$（单位向量）
\item 对应的wrench：$F_i = ([p_i] \hat{d}_i, \hat{d}_i)$
   \begin{itemize}
   \item 力矩：$m_{iz} = (p_i \times \hat{d}_i)_z$
   \item 力：$f_i = \hat{d}_i$
   \end{itemize}
\end{itemize}

\textbf{两个摩擦接触}：
\begin{itemize}
\item 接触1：有2条摩擦锥边 $\Rightarrow$ 2个wrench：$F_1, F_2$
\item 接触2：有2条摩擦锥边 $\Rightarrow$ 2个wrench：$F_3, F_4$
\item 总共：\textbf{4个wrench}
\end{itemize}

\subsection{为什么是四个Wrench？}

\textbf{数学表达}：

设两个接触：
\begin{itemize}
\item 接触1：位置$p_1$，法向量$\hat{n}_1$，摩擦系数$\mu_1$
\item 接触2：位置$p_2$，法向量$\hat{n}_2$，摩擦系数$\mu_2$
\end{itemize}

\textbf{接触1的摩擦锥边}：
\begin{itemize}
\item 边1：方向$\hat{d}_{11} = \hat{n}_1 + \mu_1 \hat{t}_1$（归一化）
\item 边2：方向$\hat{d}_{12} = \hat{n}_1 - \mu_1 \hat{t}_1$（归一化）
\item Wrench：$F_1 = ([p_1] \hat{d}_{11}, \hat{d}_{11})$，$F_2 = ([p_1] \hat{d}_{12}, \hat{d}_{12})$
\end{itemize}

\textbf{接触2的摩擦锥边}：
\begin{itemize}
\item 边1：方向$\hat{d}_{21} = \hat{n}_2 + \mu_2 \hat{t}_2$（归一化）
\item 边2：方向$\hat{d}_{22} = \hat{n}_2 - \mu_2 \hat{t}_2$（归一化）
\item Wrench：$F_3 = ([p_2] \hat{d}_{21}, \hat{d}_{21})$，$F_4 = ([p_2] \hat{d}_{22}, \hat{d}_{22})$
\end{itemize}

\textbf{总共4个wrench}：$\{F_1, F_2, F_3, F_4\}$

\section{第三部分：力闭合等价于没有一致的力矩标记}

\subsection{力闭合的定义}

\textbf{力闭合}：复合wrench锥包含整个wrench空间
\begin{equation}
\text{pos}(\{F_1, F_2, F_3, F_4\}) = \mathbb{R}^3
\end{equation}

\textbf{含义}：
\begin{itemize}
\item 可以产生\textbf{任意方向}的wrench
\item 可以平衡\textbf{任意}外部干扰
\end{itemize}

\subsection{一致力矩标记的含义}

\textbf{一致标记区域}（灰色区域）：
\begin{itemize}
\item 所有wrench对该区域内的点都产生\textbf{相同符号}的力矩
\item 例如：所有wrench都产生正力矩
\end{itemize}

\textbf{一致标记区域的问题}：
\begin{itemize}
\item 如果存在一致标记区域，说明所有wrench对某些点都产生相同符号的力矩
\item 这意味着：不能产生对这些点产生\textbf{相反符号}力矩的wrench
\item 因此：\textbf{不能正张成整个空间}
\end{itemize}

\subsection{没有一致标记的含义}

\textbf{没有一致标记区域}：
\begin{itemize}
\item 对于平面中的\textbf{任何点}，至少有一个wrench产生正力矩，至少有一个wrench产生负力矩
\item 这意味着：可以通过调整正线性组合的系数，产生\textbf{任意符号}的力矩
\item 因此：可以正张成整个空间
\end{itemize}

\subsection{等价性证明}

\textbf{定理}：力闭合 $\Leftrightarrow$ 没有一致的力矩标记

\textbf{证明思路}：

\textbf{方向1}：如果有力闭合，则没有一致标记
\begin{itemize}
\item 假设存在一致标记区域（灰色区域）
\item 则所有wrench对该区域内的点都产生相同符号的力矩
\item 这意味着不能产生相反符号的力矩
\item 因此不能正张成整个空间
\item 矛盾！因此没有一致标记区域
\end{itemize}

\textbf{方向2}：如果没有一致标记，则有力闭合
\begin{itemize}
\item 对于任何点，至少有一个wrench产生正力矩，至少有一个产生负力矩
\item 可以通过调整系数产生任意符号的力矩
\item 因此可以正张成整个空间
\item 力闭合成立
\end{itemize}

\section{第四部分："视线"（Line of Sight）的含义}

\subsection{什么是"视线"？}

\textbf{定义}：两个接触之间的\textbf{视线}是连接两个接触点的直线段。

\textbf{关键条件}：视线必须位于\textbf{两个摩擦锥内部}

\textbf{含义}：
\begin{itemize}
\item 从接触1到接触2的直线段
\item 这条线段必须完全在接触1的摩擦锥内
\item 这条线段必须完全在接触2的摩擦锥内
\end{itemize}

\subsection{为什么视线能判断力闭合？}

\textbf{几何解释}：

如果两个接触可以通过视线"看到"彼此：
\begin{itemize}
\item 说明两个摩擦锥\textbf{重叠}
\item 这意味着：可以产生从接触1指向接触2的力（在摩擦锥内）
\item 也可以产生从接触2指向接触1的力（在摩擦锥内）
\item 这两个力可以产生\textbf{任意方向}的合力
\end{itemize}

\textbf{力矩标记解释}：

如果视线在两个摩擦锥内：
\begin{itemize}
\item 接触1的摩擦锥边可以产生围绕接触2的力线
\item 接触2的摩擦锥边可以产生围绕接触1的力线
\item 四个wrench的力线可以\textbf{相互抵消}，产生任意方向的wrench
\item 因此：\textbf{没有一致标记区域}
\end{itemize}

\subsection{具体例子}

\textbf{场景}：两个接触抓取三角形物体

\textbf{接触1}：位置$p_1$，法向量$\hat{n}_1$，摩擦系数$\mu_1$
\textbf{接触2}：位置$p_2$，法向量$\hat{n}_2$，摩擦系数$\mu_2$

\textbf{检查视线}：
\begin{itemize}
\item 连接$p_1$和$p_2$的直线段
\item 检查：这条线段是否在接触1的摩擦锥内？
\item 检查：这条线段是否在接触2的摩擦锥内？
\end{itemize}

\textbf{如果视线在两个摩擦锥内}：
\begin{itemize}
\item 四个wrench（两个接触，每个2条边）可以正张成整个空间
\item 没有一致标记区域
\item \textbf{力闭合成立}
\end{itemize}

\textbf{如果视线不在摩擦锥内}：
\begin{itemize}
\item 可能存在一致标记区域
\item 不能正张成整个空间
\item \textbf{力闭合不成立}
\end{itemize}

\section{完整的理解链条}

\subsection{从接触配置到力闭合}

\textbf{步骤1}：两个摩擦接触
\begin{itemize}
\item 每个接触有摩擦锥
\item 每个摩擦锥有2条边
\item 总共：4个wrench
\end{itemize}

\textbf{步骤2}：四个wrench
\begin{itemize}
\item 平面wrench空间是3维的
\item 正张成需要$n+1 = 4$个向量
\item 四个wrench\textbf{足以}正张成三维空间
\end{itemize}

\textbf{步骤3}：力矩标记
\begin{itemize}
\item 绘制四个wrench的力矩标记
\item 检查是否有一致标记区域（灰色区域）
\end{itemize}

\textbf{步骤4}：力闭合判断
\begin{itemize}
\item 如果没有一致标记区域：力闭合
\item 如果有一致标记区域：非力闭合
\end{itemize}

\textbf{步骤5}：视线检查
\begin{itemize}
\item 如果两个接触可以通过视线"看到"彼此
\item 则没有一致标记区域
\item 因此力闭合成立
\end{itemize}

\section{图12.22(b)的解释}

\subsection{图示内容}

\textbf{图12.22(b)}显示：
\begin{itemize}
\item 两个接触（手指）抓取三角形物体
\item 两个摩擦锥（每个接触一个）
\item 连接两个接触的视线
\item 视线位于两个摩擦锥内部
\end{itemize}

\subsection{为什么这是力闭合？}

\textbf{分析}：
\begin{enumerate}
\item \textbf{四个wrench}：两个接触，每个有2条摩擦锥边
\item \textbf{视线在摩擦锥内}：说明两个摩擦锥重叠
\item \textbf{力矩标记}：四个wrench的力线可以相互抵消
\item \textbf{没有一致标记区域}：可以产生任意方向的wrench
\item \textbf{力闭合成立}
\end{enumerate}

\section{总结}

\subsection{关键理解}

\begin{enumerate}
\item \textbf{四个wrench足以正张成三维空间}
   \begin{itemize}
   \item 平面wrench空间是3维的
   \item 正张成需要$n+1 = 4$个向量
   \end{itemize}

\item \textbf{两个摩擦接触 = 四个wrench}
   \begin{itemize}
   \item 每个摩擦接触有2条摩擦锥边
   \item 每条边对应一个wrench
   \item 两个接触 = 4个wrench
   \end{itemize}

\item \textbf{力闭合 $\Leftrightarrow$ 没有一致标记}
   \begin{itemize}
   \item 一致标记区域表示wrench锥的限制
   \item 没有一致标记表示可以产生任意方向的wrench
   \end{itemize}

\item \textbf{视线判断力闭合}
   \begin{itemize}
   \item 如果两个接触可以通过视线"看到"彼此
   \item 则没有一致标记区域
   \item 因此力闭合成立
   \end{itemize}
\end{enumerate}

\subsection{实际应用}

\textbf{设计抓取时}：
\begin{itemize}
\item 确保两个接触的摩擦锥足够大
\item 确保连接两个接触的视线在摩擦锥内
\item 这样就能实现力闭合
\end{itemize}

\textbf{检查力闭合时}：
\begin{itemize}
\item 绘制力矩标记图
\item 检查是否有灰色区域（一致标记区域）
\item 或者检查视线是否在摩擦锥内
\end{itemize}

\end{document}

